\chapter*{ประมวลรายวิชา}
\addcontentsline{toc}{chapter}{ประมวลรายวิชา}

\noindent\textbf{รหัสและชื่อรายวิชา:} ฟิสิกส์สิ่งแวดล้อมและพลังงานสะอาด (Environmental Physics and Clean Energy Innovations)\\[0.2cm]
\noindent\textbf{จำนวนหน่วยกิต:} 3 (2-2-5) หน่วยกิต (ทฤษฎี 2 คาบ ปฏิบัติการ 2 คาบ ศึกษาค้นคว้าด้วยตนเอง 5 คาบต่อสัปดาห์)\\[0.2cm]
\noindent\textbf{อาจารย์ผู้รับผิดชอบรายวิชา:} ผู้ช่วยศาสตราจารย์ ดร.ชีวะ ทัศนา\\[0.2cm]
\noindent\textbf{สาขาวิชาและคณะ:} สาขาวิชาฟิสิกส์ คณะวิทยาศาสตร์และเทคโนโลยี มหาวิทยาลัยราชภัฏรำไพพรรณี

\section*{คำอธิบายรายวิชา}
ฟิสิกส์ของบรรยากาศโลก การแผ่รังสีและสมดุลพลังงาน ทัศนศาสตร์บรรยากาศ ทฤษฎีการกระเจิงแสงของอนุภาคฝุ่นละอองขนาดเล็ก (Rayleigh และ Mie Scattering) กลไกการลดทอนแสงและกฎของคอชมิเดอร์ (Koschmieder's Law) การประเมินค่าความเข้มข้นของฝุ่น PM2.5 จากภาพถ่ายสมาร์ทโฟนด้วยการเรียนรู้เชิงลึก (Deep Learning) เทอร์โมไดนามิกส์ของการเผาไหม้และฟิสิกส์เชื้อเพลิงแข็ง การผลิตถ่านอัดแท่งจากชีวมวลทางการเกษตรและการเสริมประสิทธิภาพความร้อนด้วยคาร์บอนแบล็ค กระบวนการไพโรไลซิสเพื่อการผลิตน้ำมันชีวภาพและถ่านชีวภาพ และกลไกโครงการลดก๊าซเรือนกระจกภาคสมัครใจ (T-VER) เพื่อการประเมินคาร์บอนเครดิต

\section*{วัตถุประสงค์การเรียนรู้ (Course Learning Outcomes; CLOs)}
เมื่อนักศึกษาได้ศึกษาตามหลักสูตรรายวิชานี้แล้ว จะต้องมีความรู้ความสามารถดังนี้:
\begin{enumerate}
    \item อธิบายหลักการทางฟิสิกส์บรรยากาศ ทฤษฎีการกระเจิงแสง และสมดุลพลังงานความร้อนได้อย่างถูกต้อง
    \item ประยุกต์ใช้แบบจำลองทัศนศาสตร์และโครงข่ายประสาทเทียมในการประมาณค่าฝุ่น PM2.5 จากภาพถ่ายดิจิทัลได้
    \item ดำเนินการทดสอบคุณสมบัติทางความร้อนของเชื้อเพลิงชีวมวล (ค่าความร้อน, การวิเคราะห์โดยประมาณ, การทดสอบความแข็งแรง) ตามมาตรฐาน มอก. ได้
    \item ออกแบบสูตรและกระบวนการอัดแท่งถ่านชีวมวลผสมคาร์บอนแบล็คเพื่อเพิ่มประสิทธิภาพการเผาไหม้ได้
    \item คำนวณปริมาณการลดการปล่อยก๊าซเรือนกระจกและประเมินคาร์บอนเครดิตตามระเบียบวิธี T-VER ได้
\end{enumerate}

\section*{แผนการประเมินผลการเรียนรู้}
\begin{table}[H]
\centering\small
\begin{tabularx}{\textwidth}{p{0.55\textwidth} p{0.25\textwidth} c}
\toprule
\textbf{กิจกรรมการประเมินผล} & \textbf{สัปดาห์ที่ประเมิน} & \textbf{สัดส่วนคะแนน} \\
\midrule
1. การมีส่วนร่วมในชั้นเรียนและสัมมนาวิชาการ & ตลอดภาคการศึกษา & 10\% \\
2. รายงานปฏิบัติการทดลองและการวัดค่าพลังงานความร้อน & สัปดาห์ที่ 3, 6, 9, 12 & 30\% \\
3. โครงงานวิจัยประเมิน PM2.5 ด้วย AI หรือพัฒนาเชื้อเพลิงชีวมวล & สัปดาห์ที่ 14--15 & 20\% \\
4. การสอบวัดผลกลางภาค & สัปดาห์ที่ 8 & 20\% \\
5. การสอบวัดผลปลายภาค & สัปดาห์ที่ 16 & 20\% \\
\midrule
\textbf{รวมทั้งสิ้น} & & \textbf{100\%} \\
\bottomrule
\end{tabularx}
\end{table}
